\section{Background and Related Work}

The systematic discovery of novel mathematical concepts and theorems has been transformed by recent advances in automated reasoning and machine learning. This section reviews key developments in automated mathematical discovery and establishes the theoretical foundation for our work on $$-recursive sequence spaces.

\subsection{Automated Mathematical Discovery}

The representation of mathematical relations through hypergraph structures, as introduced by \cite{halachmi2024ramanujan}, provides a fundamental framework for exploring connections between mathematical constants and functions. Their work demonstrates how vertices representing fundamental constants can be connected through edges that capture meaningful mathematical relationships, enabling systematic exploration of potential identities.

Building on this foundation, \cite{tsoukalas2025learning} developed the FERMAT system, which employs reinforcement learning to evaluate mathematical "interestingness" - a crucial metric for automated theory formation. Their approach quantifies the significance of mathematical patterns through a learned reward function that considers factors such as simplicity, generality, and surprising connections.

The integration of formal verification with discovery systems has been demonstrated by \cite{subercaseaux2023automated} in their work on discrete geometry. Their approach successfully discovered and proved new results about minimal pentagon arrangements, illustrating how automated reasoning can both generate and verify mathematical conjectures.

\subsection{Natural Language Integration and Scaling}

The NaturalProofs system developed by \cite{welleck2021naturalproofs} represents a significant advance in bridging formal mathematics with natural language understanding. Their work enables more intuitive interaction with automated theorem provers through natural mathematical language, facilitating the discovery process.

Large-scale mathematical exploration has been revolutionized by systems like AlphaEvolve, described in \cite{georgiev2025mathematical}, which combines evolutionary computation with language models to generate and evaluate mathematical hypotheses. This approach has been further validated by \cite{cao2025evaluating} in their evaluation of mathematical reasoning capabilities.

\subsection{Applications and Extensions}

Recent work has demonstrated the broad applicability of automated discovery techniques. \cite{davila2025reverie} showed how automated conjecturing systems can generate meaningful open problems in graph theory, while \cite{barket2024liouville} developed methods for generating integrable expressions.

The integration of domain-specific knowledge has been explored by \cite{servia2025physicsinformed} and \cite{liu2023learning} in their respective fields, demonstrating how automated discovery can be enhanced by incorporating field-specific constraints and prior knowledge.

\subsection{Formal Systems and Verification}

Advances in formal verification systems, as shown by \cite{xu2025isamini} and \cite{wang2025malot}, have provided robust frameworks for validating discovered results. These developments build on earlier work in autonomous learning systems \cite{angelov2018autonomous} and interpretable decision-making \cite{tschernutter2022interpretable}.

The educational implications of automated discovery tools have been explored by \cite{zarzuelo2025experience}, while \cite{wang2025polybasic} has contributed to the theoretical foundations of efficient mathematical reasoning systems.

Our work on $$-recursive sequence spaces builds upon these developments, particularly the hypergraph representation framework of \cite{halachmi2024ramanujan} and the formal verification approaches of \cite{xu2025isamini}. We extend these ideas to develop a novel mathematical structure that bridges recursive function theory and functional analysis, providing new tools for automated mathematical discovery.